\newcommand{\cert}{\textit{cert}}
\newcommand{\fals}{\textit{fals}}

\Problem{Push negations}

\Statement


We wish to trasform a Boolean expression by pushing
all the negations down to the variables. For instance,
$\lnot(a\land b)$ should be trasformed to
$(\lnot a\lor \lnot b)$.

\medskip

Expressions may contain variables (lowercase letters), conjunctions
(\texttt{'*'}), disjunctions (\texttt{'+'}) and negations (\texttt{'!'}).
For simplicity, expressions are written in their fully parenthised form.
See the examples.

\medskip

You are given an almost complete program (see attached code \texttt{code.cc}).
The main function is already written, as well as a
a simple @Formula@ class that represents Boolean expressions
using tree nodes. Here is the @Node@ structure of these trees:

\begin{lstlisting}
    struct Node {
        char op;        // operand ('a'-'z') or operator ('+' or '*')
        bool neg;       // tells if this node is negated
        Node* left;     // left subformula
        Node* right;    // right subformula
    };
\end{lstlisting}

Each leaf node contains a variable in its @op@ field.
Each non-leaf node contains \texttt{'+'} or \texttt{'*'} in its @op@ field
and pointers to its left and right children. In addition, all nodes have a @neg@
field that indicates whether or not this node is negated.

\medskip

The only attribute in the @Formula@ class is the @root@ of the tree od nodes.

\medskip

Your task is just to implement the @void push_negations()@ method of the @Formula@
class.

\medskip

Do do so, you can add private methods to the class, but you cannot alter the
@Node@ structure nor the existing methods, constructors and destructors. The
@main()@ function is already written and you do not have to modify it.

\Sample

